\chapter{INTRODUCTION}

Programming languages define the formal contract between a programmer's intent
and machine execution. Building a language processor from source text through to
evaluated output requires a coherent pipeline of components, each of which must
correctly transform its input before passing it downstream. Lexical analysis
identifies atomic tokens; parsing imposes hierarchical structure; static type
checking validates semantic correctness before execution; and interpretation
evaluates the validated program \citep{aho2006compilers,nystrom2021crafting}. When these stages are connected cleanly, errors
are caught at the earliest possible point, and the overall system remains
tractable to reason about, test, and extend.

This report presents the design and implementation of a statically-typed
interpreted language that realises this pipeline in full. The language supports
four primitive types --- Integer, Float, Boolean, and String --- and provides
arithmetic and equality expressions, assignment, conditional and iterative
control flow, user-defined functions with value parameter passing, and a
built-in callable \texttt{print()} function. Types are inferred from program context
rather than declared explicitly, combining the safety of static typing with the
conciseness of a declaration-free syntax. The implementation comprises three
principal components developed in pipeline order: the lexer and symbol table,
the recursive-descent parser with parse-tree and abstract-syntax-tree views,
and the type checker and interpreter.

The system can be invoked from the command line or through an interactive
desktop interface, both exposing five inspection views --- tokens, parse tree,
AST, inferred type table, and execution output. The parse tree supports
reporting and debugging, while the AST remains the executable intermediate
representation consumed by the type checker and interpreter. An automated
regression suite of 77 tests verifies correctness across all
stages and language features.
The report explicitly distinguishes source-level lexeme strings, token types
emitted by the lexer, and CFG terminals consumed by the parser;
Section~\ref{sec:input-to-cfg} demonstrates this separation concretely through an annotated
example.

The main contributions of this project are:

\begin{itemize}[leftmargin=1.5em, itemsep=2pt]
    \item Design of a statically-typed programming language supporting Integer,
    Float, Boolean, and String types without operator overloading.
    \item Implementation of a hand-written lexer and a recursive-descent parser
    with explicit parse-tree and abstract syntax tree generation.
    \item Development of a syntax-directed type inference and static
    type-checking system without explicit type declarations.
    \item Separation of integer and floating-point arithmetic operators
    following the Caml~Light convention to eliminate semantic ambiguity
    \citep{nystrom2021crafting}.
    \item Construction of a complete end-to-end execution pipeline consisting
    of lexical analysis, parsing, semantic analysis, and interpretation.
    \item Development of both command-line and graphical interfaces exposing
    all compiler pipeline stages for inspection and debugging.
    \item Validation of the implementation through automated regression
    testing and semantic error verification.
\end{itemize}

The remainder of this report is organised as follows. Chapter~2 presents the
language design, token specification, grammar, and operator precedence rules.
Chapters~3 and~4 cover the front end in pipeline order: lexical analysis and
the symbol table, followed by recursive-descent parsing and the AST.
Chapter~5 explains the type inference algorithm and semantic rules.
Chapter~6 presents the system architecture and implementation, including the
end-to-end pipeline, project layout, and technology stack. Chapter~7 discusses
testing and verification, and Chapter~8 concludes the report.
A supplementary walkthrough document traces the complete pipeline end-to-end
for a concrete factorial program, covering all four stages from lexical
analysis through to interpreted output.
